% ============================================================
% ABSTRACT
% ============================================================
\newpage
\thispagestyle{plain}

\begin{center}
    {\Large \textbf{Abstract}}\\[1cm]
\end{center}

This thesis presents the design and implementation of an intelligent web platform aimed at digitalizing the catering activities of \textit{Assiette Gala Sfaxienne}. The project addresses concrete operational needs, including menu catalog management, order processing, event handling, and quotation workflows.

The proposed system is built on a microservices architecture with domain-oriented services (authentication, catalog, orders, events, and AI) connected through a centralized API gateway. This architecture improves modularity, maintainability, and long-term scalability.

An AI layer is integrated to enhance user interaction through a RAG-based chatbot and a personalized recommendation engine. Development was conducted using the Scrum framework, with iterative sprint delivery and continuous functional validation.

The final outcome demonstrates the relevance of a complete, business-oriented digital solution that fits the local catering context and provides a robust foundation for future improvements.

\vspace{0.8cm}
\noindent\textbf{Keywords:} catering, microservices, artificial intelligence, RAG, recommendation system, event management, Scrum.

\clearpage
